\section{ASE}
Questa libreria \cite{ase-doc} ci permette di analizzare le simulazioni di dinamica molecolare prodotte tramite QE. Essenzialmente possiamo dire che svolge il ruolo di ponte tra la struttura fatta di coordinate e una che sia computazionalente utilizzabile. Ci permette di trasportare le informazioni che QE utilizza per svolgere i calcoli, come le condizioni di periodicità, in un formato poi accessibile anche da altre librerie.
\\

In particolare per QE abbiamo utilizzato i seguenti:

\begin{lstlisting}[language=Python]
    frame = ase.io.read(file_name, index=":", format="espresso-out")
\end{lstlisting}

dove \textit{index=":"} viene utilizzato per leggere tutti i frame presenti nel file di input e \textit{format="espresso-out"} serve per impostare la tipologia di file da leggere.
\footnote{
La specifica sezione della documentazione riguardante questo punto è reperibile al seguente \url{https://ase-lib.org/ase/io/io.html}
}

La funzione \textit{read} ci restituisce un oggetto di tipo \textit{atoms}
\footnote{
La specifica sezione della documentazione riguardante questo punto è reperibile al seguente \url{https://ase-lib.org/ase/atoms.html}
}, che costituisce l'oggetto python che contiene le informazioni che definiscono il materiale da simulare che sono state estratte dai file di output di QE.


\section{DScribe}
Questa è la libreria \cite{dscribe-doc} che si occupa di trasformare la struttura atomica definita tramite ASE in una qualcosa comprensibile dai modelli di ML. Lo scopo principale di questa procedura è quello di rendere il training di un modello il più facile possibile. L'idea è quella di creare un ogetto che identifichi in modo univoco una determinata struttura atomica, e che sia invariante per modifiche "matematiche" della struttura stessa, come possono esserlo le rotazioni; se così non fosse il modello dipenderebbe da un sistema di coordinate e andrebbe sottoposto a trining anche per le diverse rotazioni. Questo processo porta alla produzione di un descriptor, o vettore di feature. La prima cosa che facciamo notare è che separando il descriptor dalla fase di training è possibile riutilizzarlo con un diverso modello, diminuendo così il tempo necessario per l'addestramento. Abbiamo deciso di utilizzare il descriptor SOAP, per la rappresentazione locale della struttura atomica in analisi.